\section{Introduction}
\label{sec:ch3-introduction}
Smart building design, grounded in building performance engineering, primarily equates human experiences to physical ``comfort'', optimising environmental conditions of temperature, lighting, acoustics, and air quality~\cite{alsafery2023sensing, alavi2017comfort, frontczak2011literature}. Architecture, by contrast, approaches human experience through a more nuanced lens, engaging with the emotional~\cite{brady2013sublime}\footnote{\url{https://nieuweinstituut.nl/en/projects/ma-yansong-architectuur-emotie}}, social, and cultural dimensions of inhabiting space~\cite{pallasmaa2024eyes, kellert2015practice}. Rather than reducing experience to measurable parameters, architectural traditions emphasise how built environments shape perception, interpretation, and everyday use.

Recently, these domains have converged around concerns for how people experience the built environment~\cite{dourish2001action, wiberg2020interaction, kirsh2019architects} and how technologies may be designed to understand diverse human experiences~\cite{wiberg2020interaction, alavi2018artifacts}. Within human-computer interaction (HCI), this agenda is articulated through human–building interaction (HBI)~\cite{alavi2019introduction}. HBI is a branch of HCI that examines how interactive building technologies shape everyday experience in physical space~\cite{alavi2019introduction, alavi2018artifacts, alavi2017comfort}. HBI extends longstanding HCI interests like contextual interaction and experience-centred design into technology-rich architectural settings. As a result, HBI scholarship has moved beyond comfort and control, inviting inquiry into aesthetics~\cite{economidou2021moving}, social engagement~\cite{hoare2015hide}, place-making~\cite{beale2022digital}, and biophilic design~\cite{rao2025designing, rao2025lessons}.

Yet a conceptual and methodological gap remains. While the field acknowledges the experiential nature of smart buildings~\cite{wiberg2020interaction, rao2025smart}, most technologies remain oriented toward regulating environmental and occupancy data~\cite{zhong2020hilo, zeleny2024detection}. A key part of this experiential gap lies in how buildings shape and respond to affect. Affect in smart building research aligns with affective computing paradigms that model emotions as internal variables to detect and optimise~\cite{picard1999affective}. Consequently, affect is often treated as a static, quantifiable response to environmental stimuli. It is captured through physiological signals (e.g. heart rate) or behavioural cues (e.g. facial expressions)~\cite{gao2020n}, rather than as part of a deeper experiential process.  

To address this gap, we draw on the lens of Affective Interaction (AfI), which reframes affect as an outcome of ongoing interaction with systems, spaces, and others~\cite{lottridge2011affective, ahmadpour2025affective, caon2014affective}. In this view, affect is an emergent, relational, and socially situated process~\cite{lottridge2011affective, ahmadpour2025affective, caon2014affective}. We believe that an AfI perspective can capture how smart buildings participate in lived experiences by shaping how people attend to space, make meaning, and feel their way through everyday settings. 

We argue for AfI as a generative lens for understanding interactive experiences in smart buildings. Our aim is to illustrate \emph{what becomes visible} when smart building experience is approached through this lens. We present two exploratory studies in a smart university building: (1) an in-situ, open-ended survey prompting occupants to describe their experience in various spaces using an \textit{emotion} framing and a \textit{comfort} framing; and (2) a semi-structured building walk-through to trace how participants make sense of their experiences individually and collectively. We offer these studies as exemplars to demonstrate the potential of AfI to surface experiential dimensions that are largely overlooked in the current smart building research, and to inspire broader experimentation with AfI as a research perspective.

To explore this perspective, we developed two research questions. First, we examined how an AfI lens might reveal experiential dimensions that get overlooked in standard sensing or optimisation frameworks. \textbf{RQ1.} asks: \textit{What does an AfI viewpoint surface about occupant experiences in smart buildings that conventional comfort-oriented approaches tend to background or miss?} Second, we consider how AfI might reshape the goals, methods, and concepts that underpin smart building research. \textbf{RQ2.} asks: \textit{How can smart building design and evaluation methods evolve to more effectively support the complexity of inhabited experiences, surfaced through an AfI lens?}

Broadly, we contribute:
\begin{enumerate}
    \item \emph{Empirical:} Five \emph{Ways} in which occupants perceive, interpret, and adjust to their environments.
    \item \emph{Conceptual:} Six \textit{Affective Practices} (AP1---AP6) that capture recurrent strategies occupants use to regulate their experience in buildings; and corresponding \textit{Interaction Qualities} (IQ1---IQ6) for smart technologies to support these strategies.
    \item \emph{Agenda-setting:} Four directions for smart building design that include design aims, evaluation criteria, methodological focus, and epistemological commitments. 
\end{enumerate}

\textbf{HCI Contributions.} We extend experience-centred interaction research into smart buildings through an AfI lens. Conceptually, we introduce \emph{Affective Practices} and \emph{Interaction Qualities} as tools for capturing how interactive systems engage with the affective, temporal, and socially situated nature of inhabitation---dimensions under-represented in HCI–HBI accounts of smart environments. Methodologically, we complement sensor-driven approaches with situated, interpretive methods that foreground lived experience. We show how building walk-throughs and body-mapping techniques (inspired by soma-based HCI~\cite{hook2018designing}) elicit affective sense-making within everyday environments. Together, these contributions expand HCI’s methodological repertoire for examining affect as emergent and interactional in technology-rich architectural settings.



% \textbf{HCI Contributions.} We offer an experience-centred account of interaction by introducing an AfI lens for examining human experiences in smart buildings. Our \emph{Affective Practices} and \emph{Interaction Qualities} provides new conceptual tools for HCI, articulating how interactive systems can engage with affective, temporal, and socially situated dimensions of experience. Methodologically, we complement sensor-driven comfort studies with our situated methods, showing how established HCI approaches (e.g. our body maps inspired by soma research~\cite{hook2018designing}) can surface interpretive, embodied, and relational aspects of inhabiting smart buildings.}





% Specifically, we identified four overarching insights for HCI--HBI research:
% \begin{enumerate}
%     \item \emph{Affective experience in smart buildings is emergent, relational, and practice-based.} Occupants make sense of spaces through unfolding practices---somatic appraisal, anticipation, micro-adjustment, social attunement, and interpretive sense-making---rather than discrete emotional states or comfort ratings.
%     \item \emph{Experience is co-shaped by bodily sensing, spatial interpretation, and system legibility.} People continuously align their bodily state, task demands, and working theories of automation, revealing forms of interaction that remain invisible to comfort-oriented sensing or optimisation models.
%     \item \emph{Supporting experience requires designing for interpretive participation, not just environmental control.} The APs and IQs show how smart buildings can scaffold the temporal, social, and meaning-making dimensions of inhabitation, extending HCI’s interaction vocabulary into architectural settings.
%     \item \emph{Evaluating smart buildings demands situated, embodied, and temporally sensitive methods.} Our paired emotion/comfort probes, building walks, and affective body mapping demonstrate how HCI methodologies can complement sensor-driven approaches by surfacing interpretive depth and lived complexity.
% \end{enumerate}}
